
\emph{est dit} de type fini \emph{si toute famille limit\'ee de classes de faisceaux inversibles sur les fibres de \(X/S\) a comme image inverse une famille limit\'ee de classes de faisceaux sur les fibres de \(Y/S\).}

\textbf{Lemme 3.4.} \emph{Soient \(X,Y,X',Y'\) propres sur \(S\) noeth\'erien et \(T\to S\) surjectif de type fini. Alors}~:
\begin{enumerate}[label=(\roman*)]
\item \emph{Pour qu'un morphisme \(\Phi:\underline{\Pic}_{Y/S}\to\underline{\Pic}_{X/S}\) soit de type fini, il faut et il suffit que \(\Phi_T:\underline{\Pic}_{Y_T/T}\to\underline{\Pic}_{X_T/T}\) soit de type fini.}
\item \emph{Dans un diagramme commutatif}
\[
\begin{tikzcd}
\Phi':\underline{\Pic}_{Y'/S}\arrow[r]&
\underline{\Pic}_{X'/S}\\
\Phi:\underline{\Pic}_{Y/S}\arrow[r]\arrow[u,"\Phi_Y"]&
\underline{\Pic}_{X/S}\arrow[u,"\Phi_X"]
\end{tikzcd}
\]
\emph{avec \(\Phi'\) et \(\Phi_Y\) de type fini, on a \(\Phi\) de type fini.}
\end{enumerate}

En effet, supposons \(\Phi_T\) de type fini. Soit \(\mathscr L\) une famille sur les fibres de \(X/S\), limit\'ee par un \(L\) sur un \(X_R\). Ia famille \(\mathscr L_T=\mathscr L\times_S T\) sur les fibres de \(X_T/T\), d\'efinie de fa\c{c}on \'evidente, est limit\'ee par l'image inverse de \(L\) sur \(X_{R\times T}\). Si \(\Phi_T^{-1}(\mathscr L_T)\) est limit\'e par un \(M\) sur un \(Y_Q\), alors \(M\) limite aussi \(\Phi^{-1}(\mathscr L)\).

La r\'eciproque de (i) et l'assertion (ii) se d\'emontrent de fa\c{c}on analogue.

\textbf{Th\'eor\`eme 3.5.} \emph{Soient \(S\) un sch\'ema noeth\'erien, \(X\) et \(Y\) deux \(S\)-sch\'emas propres et \(f:X\to Y\) un morphisme. Si \(f\) est surjectif, alors le morphisme induit entre les foncteurs de Picard}
\[
f^*:\underline{\Pic}_{Y/S}\longrightarrow\underline{\Pic}_{X/S}
\]
